\chapter[Precipitation, Coordination, and Chemical Detection]{Precipitation,\protect\\ Coordination, and Chemical Detection}

\begin{kaichang}
Find an electric kettle used for more than six months and look inside. Its bottom and walls probably wear a gray-white crust, rough as sandpaper under your fingernail. This is limescale. A common household remedy is soaking it in vinegar: fine bubbles appear on the white crust, which gradually thins and disappears.

Do not dismiss this as unremarkable. Think carefully: tap water looks clear, and even boiled water looks clean, so \textbf{where did this white material come from}? If the water really contains hidden ingredients, why can we not see them? And vinegar does not scrub the crust, only soak it, so why does it disappear while bubbling?

Ask one step further: water utilities, mineral-water companies, and swimming-pool managers all claim to have ``tested'' water quality. How do they know what invisible, intangible substances it contains, and how much? Guess an answer. By this chapter's end, you will not only explain limescale but possess three tools of the chemical detective.
\end{kaichang}

\section{The Kettle's Invisible Guests}

First, lay out the visible facts.

Fact one: prolonged boiling of clear water deposits a white solid. Some \textbf{invisible substances} must have been dissolved there, and boiling forced them out. Fact two: vinegar attacks the solid, making it bubble and gradually dissolve. Fact three: use only bottled purified water in the same kettle, and almost no scale appears. Scale's raw material is therefore not water itself, but its passengers.

Similar phenomena abound. Kettle spouts, showerhead holes, and vacuum-flask liners acquire white, frost-like deposits over time. Caves offer grander examples: stone spires and columns hanging from ceilings are white solids delivered drop by drop, except that nature took hundreds of thousands of years. A subtler example appears when washing clothes in hard water, rich in calcium and magnesium: soap becomes ``lazy,'' producing little foam and gray-white scum floating on the surface. The invisible guests are at work again.

Who are they? To answer, zoom in to the particles.

\section{Guests Revealed: Ions Roaming in Solution}

Chapter 25 showed how much more happens in solutions than meets the eye. When salt dissolves, it does not disappear: sodium and chloride ions separate and roam, surrounded by water molecules. Chapter 23 explained how water's polarity takes ionic crystals apart. River, well, and tap water passing underground through rock and soil similarly ``invite'' ions from limestone, mainly calcium carbonate, gypsum, and other minerals. The most common guests are calcium ions, \chem{Ca^{2+}}, and bicarbonate ions, \chem{HCO_3^-}.

Ordinarily, they coexist peacefully because they move vigorously and are not highly concentrated. Boiling changes the situation: heating makes bicarbonate unstable, converting it to carbonate, \chem{CO_3^{2-}}, and releasing carbon dioxide. When carbonate meets calcium, strong electrostatic attraction, the ionic bonding of Chapter 18, joins them. Water has difficulty pulling the resulting calcium carbonate apart and carrying it away, so batches leave solution and accumulate as solid. This is limescale, mainly limestone's relative, calcium carbonate, \chem{CaCO_3}.

Stalactites and kettle scale thus perform the same drama on different stages. Water bearing calcium and bicarbonate seeps from a cave ceiling; evaporation and escaping carbon dioxide precipitate calcium carbonate bit by bit, building a column over hundreds of thousands of years.

Why does vinegar remove it? Acetic acid supplies hydrogen ions, \chem{H^+}; Chapter 32 defined acids as proton donors. Hydrogen ions capture carbonate and turn it into carbon dioxide gas and water. Those fine bubbles are carbon dioxide. The solid is dismantled into soluble ions and escaping gas, so the hard crust disappears. Precipitates can be built or dismantled: this is the beginning of our first detective tool.

\section{How Crowded Is Too Crowded? Solubility Product}

``Sparingly soluble'' is vague. How poorly does calcium carbonate dissolve? Chemists need a quantitative ruler.

Bring back Chapter 31's crucial idea: chemical equilibrium is not standing still. A calcium carbonate solid in water looks peaceful, but particles work two shifts: calcium and carbonate continually leave its surface, while dissolved ions continually strike it and recrystallize. Equal rates establish dynamic equilibrium between dissolution and precipitation, making the solution saturated.

\begin{qiaoliang}
For a sparingly soluble salt such as calcium carbonate, the product of its two ion concentrations at equilibrium is fixed at constant temperature. This is the \textbf{solubility product}, $K_{\mathrm{sp}}$:
\[K_{\mathrm{sp}}(\chem{CaCO_3}) = [\chem{Ca^{2+}}]\times[\chem{CO_3^{2-}}] \approx 3\times10^{-9}\quad(25\,^{\circ}\mathrm{C})\]
Judge it like an account: multiply current actual concentrations to obtain $Q$. Below $K_{\mathrm{sp}}$, the solution still has room and solid continues dissolving. Above $K_{\mathrm{sp}}$, it is overcrowded and excess ions must precipitate. Precipitation is not a reaction going completely to the end; this ruler has barred the excess at the door.
\end{qiaoliang}

Use the ruler for a calculation. Let calcium carbonate's solubility in pure water be $s$. Both equilibrium ion concentrations equal $s$, giving $s^2\approx3\times10^{-9}$ and $s\approx5\times10^{-5}\ \mathrm{mol/L}$. With molar mass about \ud{100}{g/mol}, that is only about \ud{5}{mg} per liter: less than a grain of salt in a cup of water. ``Sparingly soluble'' has become a number.

Now an everyday calculation. Many cities have tap-water hardness around \ud{200}{mg/L} as calcium carbonate. Suppose a \ud{1.5}{L} kettle boils one full load daily. Annual water throughput is $1.5\times365\approx550$ liters, carrying about $200\times550\approx1.1\times10^{5}\ \mathrm{mg}$ of potential scale, over a hundred grams. Even if only one-third deposits, that is over thirty grams a year, enough for a hard bottom crust. Ion concentrations and one product make an annoying household problem calculable.

Precipitation also has a medical celebrity. For gastrointestinal imaging, doctors may give patients a white ``barium meal,'' a suspension of barium sulfate, \chem{BaSO_4}. This seems contradictory: soluble barium salts are highly poisonous because barium ions interfere with nerves and muscles, yet patients drink barium meals every day. The secret is solubility product: barium sulfate has $K_{\mathrm{sp}}$ of only about $1\times10^{-10}$, dissolving roughly two milligrams per liter. Almost none is absorbed by the gut; it passes through unchanged while blocking X-rays to reveal digestive contours. Toxicity depends not merely on the element but its form and how much dissolves, a principle already encountered with oxidation states in Chapter 33.

\textbf{Translator safety note:} Low solubility alone does not establish that a barium preparation is safe to swallow. Only professionally selected medical formulations belong in an imaging procedure; aspiration and other complications are possible. Do not ingest laboratory barium sulfate or use the final exercise as medical reassurance. See the \href{https://www.accessdata.fda.gov/drugsatfda_docs/nda/2016/208844Orig1s000MedR.pdf}{FDA barium-sulfate clinical safety review}.

\section{Tool One: Make Ions Reveal Themselves}

Return to the opening question: how can we identify what is in water? The chemical detective's first tool is precipitation's \textbf{specificity}, turning a target ion into visible evidence.

The classic example tests chloride. Add silver nitrate to a water sample; a white precipitate points to chloride, because silver chloride, \chem{AgCl}, is extremely insoluble, with $K_{\mathrm{sp}}\approx1.8\times10^{-10}$, and precipitates immediately. Symbols record the invisible process in one line:
\[\chem{Ag^+} + \chem{Cl^-} \rightarrow \chem{AgCl}\downarrow\]
Similarly, barium chloride tests sulfate, \chem{SO_4^{2-}}, by producing white barium sulfate evidence. Hardness can be tested through calcium and magnesium causing scum in soapy water.

But real detectives know that \textbf{one piece of evidence does not identify the culprit}. Silver nitrate also precipitates carbonate; jumping to conclusions risks a false conviction. Standard chloride testing therefore includes a seemingly unnecessary step: first acidify with dilute nitric acid. Impostor precipitates such as silver carbonate dissolve in acid as hydrogen ions turn carbonate into escaping carbon dioxide; acid-insoluble silver chloride remains. Eliminating interference is as important to analytical chemistry as producing a signal. Detective thinking asks not ``Did it react? Then conclude,'' but ``What else could produce this observation, and how do I exclude it?''

\begin{shiyan}
\textbf{Identify Your Water with Soap}

Materials: two identical clear cups; bottled purified water; tap or mineral water; liquid soap or dishwashing detergent; a chopstick.

Procedure: half-fill one cup with purified water and the other with tap water. Add equal soap amounts, such as ten drops each. Stir with the chopstick at the same speed for the same time, or close lids tightly and shake ten times, then leave to stand and observe.

What to observe: which cup has abundant, lasting foam? Which has little foam and gray-white scum floating on it? The scum is insoluble material formed by calcium and magnesium ions combining with soap molecules. The lazier the foam, the more calcium and magnesium and the harder the water.

Safety: \textbf{do not drink} experimental water; keep soap out of your eyes and wash the cups afterward. This uses only household items and can be done with confidence.
\end{shiyan}

\section{Tool Two: Give Ions a Coat}

The first tool captures ions outside the solution. The second does the opposite: give an ion a coat to make it more conspicuous or better behaved.

Chapter 32 expanded acid--base definitions from donating protons to accepting or donating electron pairs. Follow this one step further: a metal cation is a vacant room for electron pairs, while some molecules and ions have spare pairs like eager tenants. When they approach and use these pairs to bond with the central ion, the resulting whole is a \textbf{complex ion}. Electron-pair donors are \textbf{ligands}; the cation surrounded at the center is the \textbf{central ion}; the number of directly attached donor atoms is the \textbf{coordination number}. This tenant--landlord union is a coordinate bond, a type of Chapter 19's covalent bond in which one partner supplies both electrons.

The best-known demonstration begins with pale-blue copper sulfate solution, where each copper ion is actually a complex surrounded by six water molecules, coordination number 6. A little aqueous ammonia first gives pale-blue copper hydroxide precipitate. Add excess ammonia, and the precipitate unexpectedly dissolves into a vivid deep-blue solution: ammonia displaces water to form tetraamminecopper ions, with four ammonia molecules embracing \chem{Cu^{2+}}, coordination number 4:
\[\chem{Cu^{2+}} + 4\chem{NH_3} \rightarrow [\chem{Cu(NH_3)_4}]^{2+}\]
That color change is detective evidence: the characteristic deep blue testifies to copper ions. Chapter 15's colorful transition metals owe many colors to complexes: change the central ion's coat, and its color changes.

Coordination is more than a color game. Laundry detergents often contain agents that bind calcium and magnesium, stopping hard-water ions from causing trouble. Boiler and plumbing workers use similar substances to pull scale's calcium back into solution. In medicine, one emergency treatment for heavy-metal poisoning, such as lead poisoning, is chelation therapy. An injected drug such as EDTA grips lead with six donor atoms like crab claws, making a stable, water-soluble complex excreted in urine. ``Chelation'' refers to a crab's claw: six gripping points hold more firmly than one. Such drugs can also remove needed calcium and zinc, however, so doctors must strictly control dosage and treatment duration. Never use them on your own.

Now reinterpret limescale's accounts. Calcium carbonate's low solubility is not absolute: acid removes carbonate with hydrogen ions, while complexing agents pull away calcium using electron pairs. Precipitation, dissolution, and coordination are different moves in the same equilibrium game.

\section{Tool Three: Read Colors and Weigh Molecules}

The first two tools are clever, but require guessing a suspect before choosing reagents. With an entirely unknown sample, or a target at only one part in a hundred thousand, reagent methods struggle. Modern chemical detectives use three instrumental tools with distinct roles.

First is \textbf{spectroscopy}, which makes a substance announce its name. Put a sample in a flame or arc, and excited atoms emit characteristic colors. A prism separates this light into bright lines at fixed positions, like a barcode unique to each element worldwide. Sodium's yellow lines reveal its presence, and line intensity helps estimate amount. Atomic absorption spectroscopy reverses the process: send light characteristic of an element through sample vapor, and the absorbed amount reveals how many such atoms are present. Lead, cadmium, and mercury in water and heavy metals in soil are measured this way, with sensitivity reaching parts per billion.

Second is \textbf{chromatography}, which \textbf{separates mixtures}. You may have tried this in biology: spot spinach pigment extract on filter paper and dip the strip's lower end into solvent. As solvent climbs, pigments travel at different speeds because their preferences for solvent and paper differ. They separate into green chlorophyll, yellow xanthophyll, and orange carotene bands. Modern chromatography perfects the idea: dozens of components emerge one by one from a narrow tube within minutes, each identified by its arrival time. Athlete urine testing and food pesticide-residue screening almost always begin with chromatography.

Third is \textbf{mass spectrometry}, which \textbf{weighs molecules}. Convert molecules to charged fragments, then send them through electric and magnetic fields. Heavier particles turn less readily, sorting them by mass. Measure precise molecular and fragment masses and compare with databases to identify the molecule with considerable confidence. Airport luggage swabs for explosives, police drug identification, and analysis of Martian soil all use it.

In practice, the tools often cooperate: chromatography separates the mixture, then mass spectrometry and spectroscopy identify components. This is the production line behind news that a particular substance was detected.

\begin{zhengju}
In 1860, the German chemist Bunsen and physicist Kirchhoff did something that seemed odd: evaporated several tonnes of mineral water, put the residue in a flame, and examined its light with a spectroscope. Why such effort? They had a clear question: \textbf{could elemental spectral fingerprints reveal unknown elements?}

Their design was careful. Bunsen specifically improved a high-temperature gas burner with an almost colorless flame, today's Bunsen burner, so its own color would not overwhelm the sample signal. Concentrating tonnes of water into a little residue strengthened trace-component lines enough to see. Among familiar lines in its spectrum, they found \textbf{two previously unseen bright-blue lines}.

The crucial next step was eliminating alternatives. Could known elements produce new lines under unusual conditions? They compared every known element's spectrum, finding no match. An instrumental artifact? Changing apparatus and light sources and repeating measurements left both lines fixed. Only one conclusion remained: an unknown element. They named it cesium, \chem{Cs}, from Latin for sky blue. The following year, two unfamiliar dark-red lines in an ore spectrum revealed rubidium, \chem{Rb}, named for deep red. More dramatically, astronomers found an unmatched yellow line in the solar spectrum in 1868: helium was seen on the Sun twenty-seven years before being found on Earth.

The lesson is not merely how powerful spectroscopy is, but the methodological chain: a precise question, targeted design with a colorless flame and concentrated sample, exhaustive alternative explanations through comparison with known lines, and only then the announcement of something new. Detective work is founded not on instruments but on thinking that leaves wrong conclusions nowhere to hide.
\end{zhengju}

\section{Detectives Can Fail: The Methods' Limits}

Every tool has territory outside its instruction manual.

Solubility products work well in dilute solutions but become distorted at higher concentrations. Real ions interact and ligands compete to capture them, so effective concentration, activity, differs from concentration inferred by weighing. Silver chloride is poorly soluble in pure water yet dissolves in concentrated ammonia, which complexes and removes silver ions, shifting dissolution equilibrium. Insolubility is always conditional: change temperature or introduce ligands, and recalculate.

Reagent tests depend on guessing the right suspect. They answer whether a suspected substance is present, not what everything present might be. Faced with a completely unknown sample, test kits fall silent and spectroscopy, chromatography, and mass spectrometry must enter. Instruments also have limits: every method has a \textbf{detection limit}, below which noise buries the signal. Then we can say not detected, never absolutely absent. A sensitivity gap separates those statements; do not confuse them when reading reports.

Samples themselves pose traps. Does a spoonful taken from a bag of milk powder represent the whole batch? Could the container introduce lead from outside? Detective fiction celebrates reasoning, but real laboratories most easily stumble over sampling and contamination.

\begin{wuqu}
``The news says a chemical was detected in food, so the food is poisonous and must not be eaten.'' This common panic confuses detection with harm.

More sensitive instruments make detection easier, but harm is another question. Arsenic in arsenic trioxide sounds frightening, yet traces occur in everyone's blood and every cup of seawater. Bananas naturally contain radioactive potassium-40 and remain safe to eat; even purified water can contain detectable traces of lead. The key is how far the amount lies from a potentially harmful dose. Toxicology's old saying is direct: \textbf{the dose makes the poison}. Too much water too quickly can poison you, while highly toxic botulinum toxin diluted to an exceedingly tiny fraction becomes a doctor's treatment for muscle spasms.

On seeing ``detected,'' ask not ``How could this be there?'' but: how much? What is the safety limit? Was it exceeded? Is the evidence one measurement or repeated confirmation? Risk assessment is not a yes-or-no test of presence, but arithmetic about how much, how long, and for whom. This is the attitude stressed since Chapter 12: chemical names are not frightening; ignoring dose is.
\end{wuqu}

\section{Back to the Kettle}

We can now explain the opening scene fully.

Tap water traveled through rock, dissolving calcium and bicarbonate, its invisible guests. Boiling decomposes bicarbonate into carbonate, raising the product of calcium and carbonate concentrations above calcium carbonate's solubility product. Solid precipitates layer upon layer on the bottom as scale. Vinegar's acetic acid supplies hydrogen ions that turn carbonate into carbon dioxide, those fine bubbles, and water, leaving calcium in solution. Thus the white crust is dismantled. Descaler manufacturers use the same reasoning, substituting stronger acids or specialized complexing agents for acetic acid.

You can also imagine a water laboratory's work: soap tests or complexometric titration measure hardness, using EDTA dropwise to capture calcium and magnesium and reading the total at an indicator color change. Silver nitrate methods or dedicated electrodes measure chloride; atomic absorption monitors heavy metals; chromatography and mass spectrometry screen pesticides and organic pollutants. Every compliant glass of tap water has an entire detective pipeline behind it.

\section{Volume I Ends; the Next Journey Begins}

This book's first part was called ``The Matter Detective's Toolkit.'' By Volume I's final chapter, yours is quite full. Matter consists of atoms, arranged in a periodic table by electron floors. Atoms join through ionic, covalent, and metallic bonds; intermolecular forces gather and separate molecules. Polar water dissolves everything, and solution reactions answer jointly to heat, entropy, rate, equilibrium, acids and bases, and redox. This chapter's precipitation, coordination, and instrumental analysis have brought the tools together on one investigation.

A larger mystery remains. We have asked what is in water, but have you asked what \textbf{you} are built from? Skin, hair, muscles; the rice, eggs, and beef you eat; the gases you exchange when breathing: most of these molecules are not small ions and precipitates, but intricate structures assembled from hundreds or thousands of atoms. Forensic scientists can identify someone from a hair, and mass spectrometry can weigh these molecules, but inorganic-ion rules alone cannot explain their shapes and abilities.

The answer belongs to a special atom that chains to itself, closes rings, and grows branches, assembling millions of molecules like building blocks: carbon. Volume II, How Carbon Came Alive, starts at Chapter 36 to explore why carbon is the molecular world's master framework builder, and how its frameworks build sugars, fats, proteins, and ultimately life itself. With inorganic detective skills acquired, we now investigate life.

\begin{tiaozhan}
Solubility products also predict the counterintuitive \textbf{common-ion effect}. Earlier, silver chloride's pure-water solubility was calculated as $s=\sqrt{1.8\times10^{-10}}\approx1.3\times10^{-5}\ \mathrm{mol/L}$. Now put it in \ud{0.1}{mol/L} saltwater, already rich in chloride. Let its solubility be $s'$, so $[\chem{Ag^+}]=s'$ and $[\chem{Cl^-}]\approx0.1$, since salt contributes far more chloride than dissolving silver chloride. Then
\[K_{\mathrm{sp}} = s'\times0.1 \approx 1.8\times10^{-10} \quad\Rightarrow\quad s'\approx1.8\times10^{-9}\ \mathrm{mol/L}\]
Solubility falls to roughly one seven-thousandth of the original! Intuition says a little salt should not matter, but the ledger says otherwise. This quantifies the rule that adding a substance with a common ion reduces solubility, underlying industrial salting-out recovery and more complete laboratory precipitation. Calculate the solubility if salt concentration becomes \ud{0.01}{mol/L}. You can skip this section without losing the main thread.
\end{tiaozhan}

\section*{Try It Yourself}

\begin{sikaoti}
\item Boiling hard water makes scale, whereas drying seawater leaves salt rather than scale. Use solubility products to explain why sodium chloride does not precipitate in a kettle while calcium carbonate does. Hint: compare their solubility orders of magnitude.
\item A student adds silver nitrate to a water sample, sees a white precipitate, and immediately declares chloride present. What crucial step was omitted, and which impostors does it exclude? Draw your test flowchart: sample, reagent, observation, exclusion, conclusion.
\item Draw two particle stages using circles for copper ions and small tailed balls for ammonia: a little ammonia producing precipitate, then excess ammonia dissolving it into deep-blue solution. Label central ion, ligands, and coordination number. Note that this is a human-made representation: real ions do not have these colors and shapes.
\item A barium meal contains a compound of toxic barium yet can be safely swallowed. Use solubility products to explain its safety to a worried friend, and why a toxic element and an ingestible compound are not contradictory.
\item A downstream village suspects factory pollution upstream. Investigators must ask both ``Is there lead?'' and ``What unsuspected substances are present?'' Which detective methods---reagents, spectroscopy, chromatography, mass spectrometry---suit each question? Why cannot one question rely on only one method?
\item The challenge showed silver chloride's solubility collapsing in saltwater. Why does concentrated ammonia instead increase its dissolution? Use equilibrium shifting to draw a material-flow diagram of complex formation pulling dissolution forward.
\end{sikaoti}
